%%
%% Test: mb-utils — Utility Macros — v1.1
%% Stage: utils refactor (stages 1-4)
%% 
%% Purpose: Verify that mb-utils.sty works correctly:
%%   - \latintitle and \latinauthor (with \gdef, pre-initialized)
%%   - \matin@ifcmd (with 3 arguments)
%%   - \matinbookversion and \matinbookdate (from matinbook.cls)
%% 

\documentclass{matinbook}

% Test \latintitle and \latinauthor BEFORE \begin{document}
\title{عنوان فارسی کتاب}
\author{نویسنده فارسی}
\latintitle{Latin Book Title}
\latinauthor{Latin Author Name}

\begin{document}

\maketitle

\chapter{بررسی توابع کمکی}
\label{chap:test-utils}

این سند، توابع کمکی \texttt{mb-utils} را بررسی می‌کند.

%----------------------------------------
\section{عنوان و نویسنده لاتین}
\label{sec:latin-title}

این ماکروها برای صفحات جلد دوزبانه استفاده می‌شوند.

\textbf{عنوان فارسی:} عنوان فارسی کتاب

\textbf{نویسنده فارسی:} نویسنده فارسی

\textbf{عنوان لاتین:} \makeatletter\@latintitle\makeatother

\textbf{نویسنده لاتین:} \makeatletter\@latinauthor\makeatother

%----------------------------------------
\section{بررسی نسخه}
\label{sec:version}

\textbf{نسخه MatinBook:} \matinbookversion

\textbf{تاریخ نسخه:} \matinbookdate

%----------------------------------------
\section{بررسی \texttt{\textbackslash matin@ifcmd}}
\label{sec:ifcmd}

تست \texttt{\textbackslash matin@ifcmd} با یک دستور موجود:

\makeatletter
\matin@ifcmd{chapter}{✅ دستور \texttt{chapter} تعریف شده}{❌ تعریف نشده}
\makeatother

تست \texttt{\textbackslash matin@ifcmd} با یک دستور ناموجود:

\makeatletter
\matin@ifcmd{nonexistentcommand}{✅ تعریف شده}{❌ تعریف نشده}
\makeatother

%----------------------------------------
\section{قضیه و ریاضیات}
\label{sec:theorem-math}

برای اطمینان از بارگذاری درست ماژول‌ها:

\begin{theorem}[قضیه نمونه]
    برای هر عدد حقیقی $x$، داریم $x^2 \geq 0$.
\end{theorem}

\begin{proof}
    چون مربع هر عدد حقیقی نامنفی است، نتیجه حاصل می‌شود.
\end{proof}

%----------------------------------------
\section{نتیجه‌گیری}
\label{sec:conclusion}

اگر این سند بدون خطا کامپایل شود:

\begin{itemize}
    \item ✅ \texttt{\textbackslash latintitle} و \texttt{\textbackslash latinauthor} کار می‌کنند
    \item ✅ \texttt{\textbackslash @latintitle} و \texttt{\textbackslash @latinauthor} مقدار اولیه دارند
    \item ✅ \texttt{\textbackslash matinbookversion} از \texttt{matinbook.cls} می‌آید
    \item ✅ \texttt{\textbackslash matin@ifcmd} با ۳ آرگومان کار می‌کند
    \item ✅ محیط‌های علمی کار می‌کنند
\end{itemize}

\textbf{وضعیت تست:} قبول

\end{document}
